\chapter{Каналы свидетельств и калибровка уверенности}
\label{ch:evidence}

\section{Почему одного скаляра мало, но отказываться от него рано}

В базовом тексте уверенность уже введена как знаковый скаляр, и для Stage~1 это
правильная дисциплина. Но в курсе полезно увидеть, из чего этот скаляр \emph{собирается}.
Иначе он начинает казаться либо магической меткой, либо скрытым синонимом истины.

\begin{principle}[Скаляр остаётся интерфейсом, разложение живёт внутри]
Внешнее поле $\conf(v)$ удобно держать одномерным, чтобы не раздувать каждый
оператор памяти. Но внутренняя логика его обновления может и должна различать
каналы свидетельств.
\end{principle}

\section{Лог-отношение правдоподобия}

Удобно перейти от скаляра уверенности к скрытой переменной лог-правдоподобия
\[
L(v)=\log\frac{P(v\ \text{истинно})}{P(v\ \text{ложно})}.
\]
Тогда
\[
P(v\ \text{истинно})=\frac{1}{1+e^{-L(v)}},
\qquad
\conf(v)=2P-1=\tanh\frac{L(v)}{2}.
\]

Это не меняет внешнюю семантику, но делает накопление свидетельств аддитивным:
независимые вклады можно просто суммировать.

\section{Каналы свидетельств}

В расширенной версии вводится следующий внутренний разбор:
\[
L(v)=b_{\mathrm{type}}(v)+
\sum_i \rho_i\, s_i(v)-\Delta_{\mathrm{dep}}(v)-\Delta_{\mathrm{stale}}(v).
\]

Здесь:
\begin{itemize}
\item $b_{\mathrm{type}}(v)$ — prior по типу узла;
\item $s_i(v)$ — сила свидетельства канала $i$;
\item $\rho_i$ — калибровочный коэффициент надёжности канала;
\item $\Delta_{\mathrm{dep}}$ — штраф за зависимость источников;
\item $\Delta_{\mathrm{stale}}$ — штраф за устаревание.
\end{itemize}

\subsection*{Разумный минимальный набор каналов}
\begin{enumerate}
\item \textbf{Наблюдение}: узел возник как прямой observation.
\item \textbf{Внешний отчёт}: информация пришла из текста, речи, базы знаний.
\item \textbf{Вывод}: утверждение получено fusion или правилом.
\item \textbf{Само-модель}: утверждение относится к состоянию и свойствам самой Рин.
\item \textbf{Социальное эхо}: повторение информации многими зависимыми источниками.
\end{enumerate}

\begin{pitfall}[Повтор не равен независимости]
Пять пересказов одной и той же новости — не пять независимых свидетельств. Если не
штрафовать зависимость источников, социальный шум начнёт раздувать уверенность как
настоящий многократный опыт.
\end{pitfall}

\section{Штраф за зависимость}

Пусть утверждение $v$ поддержано набором источников $J(v)$. Если источники заведомо
не независимы, вместо наивной суммы полезно использовать
\[
\Delta_{\mathrm{dep}}(v)=
\lambda \sum_{i<j}\mathrm{overlap}(i,j),
\]
где $\mathrm{overlap}(i,j)$ оценивает общность происхождения: общий документ,
одинаковый provenance-root, одинаковую цепочку ретрансляции.

Практический смысл прост: память не должна принимать видимость консенсуса за
доказательство.

\section{Устаревание и темпоральная валидность}

Не всякое старение — забывание. Для некоторых узлов старение уменьшает \emph{доступность},
но не меняет истинность. Для других — меняет именно истинностный статус.

\begin{definition}[Эпистемическое старение]
Изменение $L(v)$ вследствие того, что утверждение могло утратить актуальность во
времени, даже если память о нём сохранилась.
\end{definition}

Отсюда
\[
\Delta_{\mathrm{stale}}(v)=g(\text{age}(v),\text{node\_kind}(v),\text{validity\_window}(v)).
\]
Для биографического факта функция может быть почти нулевой, для текущего места
жительства — существенной, для настроения — очень быстрой.

\section{Пороги принятия и зона честного незнания}

Скрытый логит полезен ещё и тем, что пороги легче задавать на нём, чем на интуитивной
шкале «уверена/не уверена». Но наружу система по-прежнему может говорить на языке
$\conf$.

Вводятся три зоны:
\[
|\,\conf(v)\,|<h \Rightarrow \text{честное незнание},
\]
\[
\conf(v)\ge \theta_+ \Rightarrow \text{рабочее принятие},
\qquad
\conf(v)\le \theta_- \Rightarrow \text{рабочее отвержение}.
\]

Промежуточная область между $h$ и $\theta_+$ особенно важна для Рин как существа,
которое должно \emph{не только помнить, но и осторожно формулировать}. Именно здесь
возникает реплика вида: «это похоже на X, но я бы проверила Y».

\section{Новая добавка: режим evidence-aware revision}

Это ещё одно расширение сверх базового текста. В исходном roadmap ревизия уже есть,
но теперь она делается более расчленённой.

Пусть пришло контрсвидетельство $c$ к утверждению $v$. Тогда ревизия может идти по
трём уровням.
\begin{enumerate}
\item \textbf{Локальное понижение}: пересчитать $L(v)$ без перестройки соседей.
\item \textbf{Структурный пересмотр}: пометить потомков как \texttt{need\_to\_revise}.
\item \textbf{Канальный аудит}: проверить, какой именно канал раньше доминировал и не был
ли он систематически переоценён.
\end{enumerate}

Третий пункт принципиален. Он превращает отдельную ошибку памяти в материал для
мета-обучения: если снова и снова проваливается один и тот же канал, проблема уже
не в конкретном узле, а в калибровке $\rho_i$.

\section{Worked example: слух, наблюдение и ревизия}

\begin{workedbox}[Как один и тот же факт может менять статус]
Пусть утверждение $v$ таково: «у человека $Z$ сменилось место работы».

Шаг 1. Приходит внешний отчёт из одного источника. Возникает канал \emph{reported},
уверенность умеренная.

Шаг 2. Позднее Рин видит новую публичную страницу профиля. Добавляется канал
\emph{observed}, логит растёт сильнее, чем от второго пересказа того же слуха.

Шаг 3. Затем обнаруживается, что оба ранних текста ссылаются на один и тот же
агрегатор. Возрастает $\Delta_{\mathrm{dep}}$, прежний «консенсус» частично сдувается.

Шаг 4. Через месяц находится противоречащая запись с более свежей датой. Срабатывает
темпоральный штраф к старому утверждению и запускается revision cascade для всех
выводов, где старое место работы использовалось как посылка.
\end{workedbox}

\section{Зачем это всё для Stage~1}

Можно возразить, что такая схема слишком богата для первого работающего слоя. Это
возражение верно, если понимать её как требование немедленной полной реализации.
Но неверно, если понимать её как \emph{правильный язык проектирования}.

Для Stage~1 достаточно:
\begin{enumerate}
\item хранить единый $\conf(v)$;
\item помнить \texttt{node\_kind};
\item различать хотя бы observation/report/inference;
\item логировать provenance так, чтобы зависимость источников потом было к чему привязать.
\end{enumerate}

То есть расширенная математика не отменяет первого этапа. Она задаёт рамку, в
которой первый этап не превращается в тупик.
